\documentclass[10pt,a4paper]{amsart}
\usepackage[margin=19mm]{geometry}
\usepackage{amsmath,amssymb,mathtools}
\usepackage[scr=rsfs,cal=euler]{mathalfa}
\usepackage{xfrac}
\usepackage[unicode,hidelinks,bookmarks=false]{hyperref}
\newcommand{\ie}{i.e.\ }
\newcommand{\cf}{cf.\ }
\newcommand{\resp}{respectively\ }
\newcommand{\Subsection}{\subsection}
\providecommand{\nobreakdash}{\nobreak\hbox{-}}
\setlength{\emergencystretch}{2em}
\multlinegap0pt
\usepackage{amssymb}
\usepackage{xcolor}
\usepackage{enumerate}
\newtheorem{prop}{Proposition}
\newcommand{\dr}{\mathrm{d}}
\newcommand{\ii}{\mathrm{int}}
\newcommand{\iii}{\mathrm{i}}
\newcommand{\mr}{\mathrm}
\newcommand{\ra}{\rightarrow}
\newcommand{\wh}{\widehat}
\newcommand{\xra}{\xrightarrow}
\title{BV pushforward as a quasi-isomorphism}
\author{Alberto S. Cattaneo, Pavel Mnev}
\date{Source: arXiv:2605.30558v2 (CC BY 4.0)}
\begin{document}
\maketitle

\noindent\emph{Source and licence.} BV pushforward as a quasi-isomorphism. Version arXiv:2605.30558v2 (CC BY 4.0), distributed by arXiv under the Creative Commons Attribution 4.0 International licence, http://creativecommons.org/licenses/by/4.0/, as recorded in the arXiv licence metadata for this exact version and shown on the abstract page https://arxiv.org/abs/2605.30558v2. Full licence notice: \texttt{licenses/arxiv-2605.30558-CC-BY-4.0.txt}.\par\smallskip

\noindent\textit{Editorial scope.} Counted unit: the article's prop loop (source0.tex lines 1640--1645). The article's complete original proof of it is delivered with it (source0.tex lines 1646--1673, one \texttt{proof} environment of 28 lines). No mathematics is written, repaired or re-derived: the statement and the proof are the article's own lines, in the article's own notation, and no retained line is substituted or re-flowed.  Retained uncounted as the source-local context the proof needs: lines 657--659.  Nothing was omitted from any retained span, so the package carries no omission record.  No retained line contains a citation, so the wrapper carries no bibliography and no bibliography entry.  No retained line contains a figure environment, an image-inclusion command or drawing code, so no image is delivered and the package declares no asset dependency.  Editorial formatting only, in addition: an AMS \texttt{amsart} preamble that restates the article's own declarations and macros used by the retained lines, namely the environments \texttt{prop} on separate counters, together with the standard packages those lines need.  The retained environments print in this document as Proposition 1 in order of appearance, whereas the article prints the same items with the numbering of its own full document.  The complete native source of the article as distributed in the arXiv e-print for this exact version is bundled unchanged as \texttt{sources/201681/source0.tex}.
\begin{equation}\label{Ksym=nu whK}
    K^\mr{sym}=\nu \wh{K}.
\end{equation}

\begin{prop}\label{prop: lifting ab Wilson loop}
    The lifting of $W_\gamma$ to abelian $BF+B^2$ theory is given by
    \begin{equation}\label{ab BF+BB: i_int(W)}
        i_\ii(W_\gamma)=\exp \left(\iii\alpha \oint_\gamma (A+\lambda 2 G \dr^*B_-)\right).
    \end{equation}
\end{prop}

\begin{proof}
%Let us compute $i_\ii(W_\gamma)$ -- the lifting of $W_\gamma$ to the abelian $BF+B^2$ theory. 
By homological perturbation lemma, we have
\begin{equation}\label{ab BF+BB: HPT series}
    i_\ii(W_\gamma)=(i-K(Q_\ii-\iii\hbar \Delta) i + K(Q_\ii-\iii\hbar \Delta)K(Q_\ii-\iii\hbar \Delta) i-\cdots) W_\gamma
\end{equation}
with 
\begin{multline}
    Q_\ii=\{S_\ii,-\}= \int_M (P_+ \dr\tau) \frac{\delta}{\delta B_+}+ (P_- \dr A) \frac{\delta}{\delta B_-^+}+\dr B_- \frac{\delta}{\delta A_+} + \dr B_+^+ \frac{\delta}{\delta \tau^+}\\
    +\lambda \left(
    \phi\frac{\delta}{\delta c}+\tau \frac{\delta}{\delta A}+B_+\frac{\delta}{\delta B_+^+}+B_-\frac{\delta}{\delta B_-^+}+A^+\frac{\delta}{\delta \tau^+}+c^+\frac{\delta}{\delta \phi^+}
    \right).
\end{multline}
Next we compute the terms in the series (\ref{ab BF+BB: HPT series}):
\begin{multline}
    W_\gamma(A) \xra{Q_\ii-\iii\hbar\Delta} \iii\alpha\left(-\oint_\gamma \lambda \tau\right)W_\gamma(A) \xra{K} \iii\alpha\left(-\oint_\gamma \lambda 2G\dr^* B_-\right)W_\gamma(A) \\
    \xra{Q_\ii-\iii\hbar\Delta} (\iii\alpha)^2\left(-\oint_\gamma \lambda 2G\dr^* B_-\right)\left(-\oint_\gamma \lambda \tau\right)W_\gamma(A) 
    \xra{K} 
    \frac12  (\iii\alpha)^2\left(-\oint_\gamma \lambda 2G\dr^* B_-\right)^2 W_\gamma(A)\\
    \cdots \ra (K(Q_\ii-\iii\hbar\Delta))^k (W_\gamma)=
    \frac{1}{k!}(\iii\alpha)^k \left(-\oint_\gamma \lambda 2G\dr^* B_-\right)^k W_\gamma(A).
\end{multline}
Here we use the formula (\ref{Ksym=nu whK}) for $K$, $K=\nu\wh\kappa$, and the operator $\nu$ is responsible for the appearance of the factor $\frac{1}{k!}$ above. Thus, we have
\begin{multline}
    i_\ii(W_\gamma)= \sum_{k\geq 0}\frac{(\iii\alpha)^k}{k!} \left(\oint_\gamma \lambda 2G\dr^* B_-\right)^k W_\gamma(A)\\
    =\exp \left(\iii\alpha \oint_\gamma (A+\lambda 2 G \dr^*B_-)\right).
\end{multline}
\end{proof}

\end{document}
